\documentclass{article} % For LaTeX2e
\usepackage[T1]{fontenc}
\usepackage{iclr2027_conference,times}

% Optional math commands from https://github.com/goodfeli/dlbook_notation.
\input{math_commands.tex}

\usepackage{hyperref}
\usepackage{url}
\usepackage{amsthm}
\usepackage{booktabs}
\usepackage{graphicx}

\theoremstyle{plain}
\newtheorem{theorem}{Theorem}[section]
\newtheorem{lemma}[theorem]{Lemma}
\newtheorem{proposition}[theorem]{Proposition}
\newtheorem{corollary}[theorem]{Corollary}
\newtheorem{conjecture}[theorem]{Conjecture}
\theoremstyle{definition}
\newtheorem{definition}[theorem]{Definition}
\newtheorem{assumption}[theorem]{Assumption}
\theoremstyle{remark}
\newtheorem{remark}[theorem]{Remark}

% notation
\newcommand{\Zp}{\mathbb{Z}_p}
\newcommand{\Fp}{F}                       % positive frequencies
\newcommand{\marg}{\operatorname{margin}}
\newcommand{\supp}{\operatorname{supp}}
\newcommand{\Gam}{\Gamma}
\newcommand{\cov}[2]{g_{#1}(#2)}          % coverage kernel
\newcommand{\clean}{g}                    % clean part
\newcommand{\gam}{\gamma_q}
\newcommand{\Pq}{P_q}
\newcommand{\Kstar}{K^{*}}

\title{Explaining Feature Sparsity Through the\\ Margin--Regularization Trade-off}

% Authors must not appear in the submitted version.
\author{Anonymous authors}

\newcommand{\fix}{\marginpar{FIX}}
\newcommand{\new}{\marginpar{NEW}}

%\iclrfinalcopy % Uncomment for camera-ready version, but NOT for submission.
\begin{document}

\maketitle

\begin{abstract}
Explaining the number of features in a network has been a long-desired result in mechanistic interpretability. \citet{morwani2024feature} laid the baseline with their principle of margin-maximization, showing provably that in networks where the regularization ($L_{2,\nu}$) matches its homogeneity degree ($\nu$), all features would be used for modular addition or group composition tasks. Empirically, however, networks trained with ordinary weight decay consistently show a sparse selection of features. This paper investigates why.
Our inquiry focuses specifically on modular addition and group composition tasks in polynomial and multilinear models. Our primary theoretical finding centers on the cost of new features---a trade-off between the logit margin of the solution and regularization. While more features improve margin, they cost the model in regularization. Hence we show that an optimal number of features can be calculated for networks based on the task, their homogeneity degree and the regularization. Further, our theory explains the transition in the number of features from dense to sparse---as seen empirically in \citet{mccracken2025}---by showing that this transition arises on changing either the homogeneity degree or regularization of the network.
Specifically, we empirically show that the normalized margin attained by these networks peaks at the predicted optimum for the network, matching the result of \citet{morwani2024feature} for their specific case. The optimums are also consistent with the empirical results seen in \citet{chughtai2023toy}.
\end{abstract}

\section{Introduction}
\label{sec:intro}
% TODO (user): framing paragraphs. Placeholder outline only.
\begin{itemize}
  \item Networks trained on modular addition learn a small set of Fourier features \citep{nanda2023progress,chughtai2023toy}; under the $L_{2,\nu}$ penalty the max-margin solution is provably dense \citep{morwani2024feature}; under plain weight decay trained networks are sparse, and the count moves from dense to sparse as the degree changes \citep{mccracken2025}.
  \item Question: what decides the number of features? Answer: the trade-off between coverage (margin) and the price of a feature under the regulariser, summarised by one exponent $s=\nu/q$.
  \item Contributions: (i) a universal margin bound for any neuron-additive homogeneous model; (ii) the exact optimum and its feature count in the interference-free class, with every feature carried by exactly $\nu$ neurons at root-of-unity phases; (iii) the same programme on finite groups; (iv) experiments: how interference enters trained networks, and the count under the $\nu$ and $q$ dials.
\end{itemize}

\input{sections/theory}

\section{Group composition}
\label{sec:groups}
% TODO: brief from s5theory/S5_THEORY.md §10. Skeleton only.
For group composition the same argument gives an exact margin ceiling in which each irreducible representation pays a price set by how cheaply the matrix-multiplication tensor can be written as a sum of products (at even degree the schoolbook algorithm, $d^\nu$ neurons per $d$-dimensional irrep). Among networks without same-irrep interference this ceiling is attained and fixes the feature count: one irrep on $S_5$, the Galois pair on $A_5$, as observed. Unlike on cyclic groups of prime order, networks with structured interference can exceed it on some groups ($D_3$, $S_4$, $\mathbb{Z}_{12}$), so there the count is exact only for the interference-free class.

\section{Experiments}
\label{sec:experiments}
% TODO: four items in the user's order.
\subsection{Where interference enters trained networks}
% clean fraction vs q (RESULTS_Q): 0.4 at q=2 -> 0.99 at q=nu; local-vs-global objective (LADDER §5.2).
\subsection{The exponent $s=\nu/q$}
% q sweep at fixed nu=3: count 5 -> dense, monotone; endpoint exact to 3e-6.
\subsection{Moving the degree and the penalty separately}
% counts vs nu at q=2; counts vs q at nu=3,4; matched-s triple.

\section{Discussion}
\label{sec:discussion}
% TODO

\subsection*{Reproducibility statement}
All theorems are stated with their hypotheses in \S\ref{sec:theory} and proved in full in Appendix~\ref{app:proofs}; the constants of every construction were checked numerically to machine precision. Experimental settings are listed in \S\ref{sec:experiments}; code and configuration files are provided as supplementary material.

\bibliography{refs}
\bibliographystyle{iclr2027_conference}

\appendix
\input{sections/proofs}

\end{document}
